\section*{Chapter \ref{ch:g10:why-clouds-form} --- Why Clouds Form}

\begin{solution}{exo:g10:why-clouds-form:1}
The \omterm{def:g10:why-clouds-form:humidity}{relative humidity} is the air's content of water vapour divided by its
\omterm{def:g10:why-clouds-form:humidity}{saturation} content at the same temperature and pressure, as a percentage.
The \omterm{def:g10:why-clouds-form:dew-point}{dew point} is the temperature to which the air must be cooled, keeping
its water content and pressure, to become \omterm{def:g10:why-clouds-form:humidity}{saturated}.
\end{solution}

\begin{solution}{exo:g10:why-clouds-form:2}
\num{3.8}, \num{7.6}, \num{14.7} and \qty{27.2}{g/kg}. Each
\qty{10}{\degreeCelsius} multiplies it by about 2 ($7.6/3.8 = 2.0$,
$14.7/7.6 \approx 1.9$, $27.2/14.7 \approx 1.9$).
\end{solution}

\begin{solution}{exo:g10:why-clouds-form:3}
\omterm{def:g10:why-clouds-form:humidity}{Saturation} content at \qty{24}{\degreeCelsius}: \qty{18.9}{g/kg}; \omterm{def:g10:why-clouds-form:humidity}{relative
humidity} $10.0/18.9 \approx \qty{53}{\percent}$. The \omterm{def:g10:why-clouds-form:humidity}{saturation} content is
\qty{10.0}{g/kg} at \qty{14}{\degreeCelsius}: the \omterm{def:g10:why-clouds-form:dew-point}{dew point}.
\end{solution}

\begin{solution}{exo:g10:why-clouds-form:4}
In summer the air is humid and its \omterm{def:g10:why-clouds-form:dew-point}{dew point} is well above the temperature
of the glass: the air touching the glass is cooled below its \omterm{def:g10:why-clouds-form:dew-point}{dew point} and
water condenses on it. Very dry winter air indoors has a \omterm{def:g10:why-clouds-form:dew-point}{dew point} lower
than the temperature of the glass: cooled against it, the air is still not
\omterm{def:g10:why-clouds-form:humidity}{saturated}.
\end{solution}

\begin{solution}{exo:g10:why-clouds-form:5}
Heating from below (afternoon cumulus over sunny ground); mountains (wind
forced up a range); fronts (warm air rising over cold air); convergence
(winds flowing together near the equator or into a depression).
\end{solution}

\begin{solution}{exo:g10:why-clouds-form:6}
$125 \times (30 - 18) = \qty{1500}{m}$; $125 \times (22 - 20) =
\qty{250}{m}$. The second day: clouds form only \qty{250}{m} above the
ground.
\end{solution}

\begin{solution}{exo:g10:why-clouds-form:7}
At \qty{1}{km}: $18 - 9.8 = \qty{8.2}{\degreeCelsius}$; at \qty{2}{km}:
$18 - 19.6 = \qty{-1.6}{\degreeCelsius}$. Sinking, it warms at the same rate
and reaches the ground at \qty{18}{\degreeCelsius} again.
\end{solution}

\begin{solution}{exo:g10:why-clouds-form:8}
With \qty{4}{\degreeCelsius/km}, both an unsaturated parcel
(\qty{9.8}{\degreeCelsius/km}) and a \omterm{def:g10:why-clouds-form:humidity}{saturated} one
(\qty{5}{\degreeCelsius/km}) cool faster than their surroundings: \omterm{def:g10:why-clouds-form:stability}{stable} in
both cases. With \qty{12}{\degreeCelsius/km}, even an unsaturated bubble
cools more slowly than its surroundings: it stays warmer and keeps rising,
and the bubbles of warm air rise freely.
\end{solution}

\begin{solution}{exo:g10:why-clouds-form:9}
$14.7 - 7.6 = \qty{7.1}{g}$ per kilogram of air. The \omterm{def:g10:earth-energy-budget:latent-heat}{latent heat} warms the
air in which the water condenses, slowing its cooling.
\end{solution}

\begin{solution}{exo:g10:why-clouds-form:10}
Volumes grow as the cube of the diameter: $(0.5/0.02)^3 = 25^3 = 15\,625$
droplets for a drizzle drop; $(2/0.02)^3 = 100^3 = 1\,000\,000$ for a
raindrop.
\end{solution}

\begin{solution}{exo:g10:why-clouds-form:11}
Under a clear sky the ground radiates infrared and cools, and so does the
air touching it. This cold air is denser than the air above; on slopes it
flows downhill and collects in the valley bottom. There it keeps cooling
through the night until it reaches its \omterm{def:g10:why-clouds-form:dew-point}{dew point}, and droplets form in the
air itself: fog. On the hilltops the air is not as cold and stays above its
\omterm{def:g10:why-clouds-form:dew-point}{dew point}.
\end{solution}

\begin{solution}{exo:g10:why-clouds-form:12}
At height $z$ the temperature is $T - a z$ and the \omterm{def:g10:why-clouds-form:dew-point}{dew point} $T_d - b z$;
they are equal when $(a - b) z = T - T_d$, so $z = (T - T_d)/(a - b)$. With
$b = 1.7$: $a - b = 8.1$, and $1000/8.1 \approx \qty{123}{m}$ per degree.
\end{solution}

\begin{solution}{exo:g10:why-clouds-form:13}
\omterm{def:g10:why-clouds-form:condensation-level}{Condensation level}: $125 \times (15 - 11) = \qty{500}{m}$; temperature there:
$15 - 9.8 \times 0.5 \approx \qty{10.1}{\degreeCelsius}$. Crest: $10.1 - 5
\times 1.5 = \qty{2.6}{\degreeCelsius}$. Down \qty{2}{km} at the dry rate:
$2.6 + 9.8 \times 2 \approx \qty{22.2}{\degreeCelsius}$, about seven degrees
warmer than on the windward side.
\end{solution}

\begin{solution}{exo:g10:why-clouds-form:14}
At \qty{-20}{\degreeCelsius}: $126/103 \approx 1.22$, \qty{22}{\percent}
oversaturated for the ice; at \qty{-10}{\degreeCelsius}: $287/260 \approx
1.10$, \qty{10}{\percent}. Among these temperatures, the colder the cloud,
the larger the advantage of the ice: the process works best in the cold
upper part of a cloud where droplets and crystals are mixed.
\end{solution}

\begin{solution}{exo:g10:why-clouds-form:15}
Unsaturated: the parcel cools by \qty{9.8}{\degreeCelsius}, its surroundings
by \qty{6.5}{\degreeCelsius}: it is \qty{3.3}{\degreeCelsius} colder and
sinks back. \omterm{def:g10:why-clouds-form:humidity}{Saturated}: it cools by \qty{5}{\degreeCelsius}, so it is
\qty{1.5}{\degreeCelsius} warmer and keeps rising. The air is \omterm{def:g10:why-clouds-form:stability}{unstable} only
on one condition: that the rising air be \omterm{def:g10:why-clouds-form:humidity}{saturated}.
\end{solution}

\begin{solution}{pb:g10:why-clouds-form:1}
\textbf{1.} \qty{14.7}{g/kg}.

\textbf{2.} $11.4/14.7 \approx \qty{78}{\percent}$.

\textbf{3.} \qty{16}{\degreeCelsius}, where the \omterm{def:g10:why-clouds-form:humidity}{saturation} content is
\qty{11.4}{g/kg}.

\textbf{4.} $11.4/24.1 \approx \qty{47}{\percent}$.

\textbf{5.} No: it holds the same \qty{11.4}{g/kg}. Its \omterm{def:g10:why-clouds-form:humidity}{relative humidity}
fell because, warmer, it could hold more.

\textbf{6.} At its \omterm{def:g10:why-clouds-form:dew-point}{dew point}, \qty{16}{\degreeCelsius}.

\textbf{7.} $11.4 - 8.7 = \qty{2.7}{g}$ per kilogram of air.

\textbf{8.} Under a clear night sky, grass and roofs radiate infrared and
become colder than the air; the air touching them is cooled below its \omterm{def:g10:why-clouds-form:dew-point}{dew
point} before the rest of the air is.

\textbf{9.} $(28.0 - 13.4)/1.5 \approx \qty{9.7}{\degreeCelsius/km}$,
practically the dry rate: the afternoon bubbles have mixed this layer.

\textbf{10.} $(13.4 - 0.4)/2.0 = \qty{6.5}{\degreeCelsius/km}$.

\textbf{11.} \omterm{def:g10:why-clouds-form:stability}{Stable} for an unsaturated parcel, which cools faster
(\qty{9.8}{\degreeCelsius/km}); \omterm{def:g10:why-clouds-form:stability}{unstable} for a \omterm{def:g10:why-clouds-form:humidity}{saturated} one, which cools
more slowly (\qty{5}{\degreeCelsius/km}).

\textbf{12.} The temperature rises with height, from \num{-0.3} to
\qty{1.2}{\degreeCelsius}: a very \omterm{def:g10:why-clouds-form:stability}{stable} layer, which acts as a lid. The
rising cloud stops there: its tops are near \qty{4}{km}.

\textbf{13.} $28 - 16 = \qty{12}{\degreeCelsius}$.

\textbf{14.} Temperature $28 - 9.8 = \qty{18.2}{\degreeCelsius}$; \omterm{def:g10:why-clouds-form:dew-point}{dew point}
$16 - 1.8 = \qty{14.2}{\degreeCelsius}$.

\textbf{15.} From 12 to $18.2 - 14.2 = 4.0$: it shrank by
\qty{8}{\degreeCelsius}, a rate of \qty{8}{\degreeCelsius} per kilometre.

\textbf{16.} One degree of difference is closed in $1/8$ km, that is
\qty{125}{m}.

\textbf{17.} $125 \times (20 - 16) = \qty{500}{m}$.

\textbf{18.} The \omterm{def:g10:why-clouds-form:dew-point}{dew point} stays at \qty{16}{\degreeCelsius} while the air
near the ground warms: the difference grows, and so does the height at which
it is closed.

\textbf{19.} $28 - 9.8\,z = 16 - 1.8\,z$.

\textbf{20.} $8\,z = 12$, so $z = 1.5$: the base of the afternoon cumulus is
\textbf{about \qty{1.5}{km}} above the plain (the exact computation with the
properties of water gives \qty{1.51}{km}).
\end{solution}
